\documentclass{article}
\usepackage[T1]{fontenc}
\usepackage[utf8]{inputenc}
\usepackage{lmodern}
\usepackage{amsmath}
\usepackage{sansmath}
\AtBeginDocument{\sansmath}
\renewcommand{\familydefault}{\sfdefault}
\usepackage[paperwidth=19cm,paperheight=14cm,margin=0.4cm]{geometry}
\usepackage{graphicx}
\usepackage{tikz}
\usetikzlibrary{positioning,fit,calc,arrows.meta,backgrounds,decorations.pathreplacing}
\pagestyle{empty}
\setlength{\parindent}{0pt}
\graphicspath{{img/}}

\definecolor{ink}{HTML}{0A0A0A}
\definecolor{muted}{HTML}{6A6A69}
\definecolor{line}{HTML}{BDBDBA}
\definecolor{red}{HTML}{E03A3A}
\definecolor{blue}{HTML}{2E6FE0}
\definecolor{yellow}{HTML}{D9A11A}
\definecolor{green}{HTML}{23A06A}
\definecolor{model}{HTML}{8C8C8A}

\tikzset{
  every picture/.style={font=\sffamily\small, line cap=round, line join=round},
  flecha/.style={-{Stealth[length=2.3mm,width=1.7mm]}, line width=0.7pt, draw=ink},
  grupo/.style={draw=#1, dash pattern=on 3pt off 2pt, line width=0.8pt, rounded corners=2pt, inner sep=6pt},
  titulo/.style={font=\sffamily\small\bfseries, text=#1, anchor=south west, inner sep=1pt},
  foto/.style={inner sep=0pt, outer sep=0pt, draw=ink, line width=0.5pt},
  caja/.style={draw=ink, line width=0.5pt, fill=#1, text=white, font=\sffamily\scriptsize\bfseries,
    text width=3.75cm, minimum height=0.46cm, align=center, inner sep=1.5pt},
  nota/.style={font=\sffamily\scriptsize, text=muted, align=center, inner sep=1pt},
  etapa/.style={font=\sffamily\footnotesize\bfseries, text=ink, align=center, inner sep=1pt},
  paso/.style={draw=ink, line width=0.5pt, fill=#1, text=ink, font=\sffamily\scriptsize, align=center,
    inner sep=3pt, minimum height=0.95cm, text width=2.35cm},
  paso/.default=white,
  grupofondo/.style={draw=#1, line width=0.8pt, fill=#1!6, rounded corners=3pt, inner sep=7pt},
  llave/.style={decorate, decoration={brace, mirror, amplitude=5pt, raise=2pt}, line width=0.7pt, draw=#1},
  llavetexto/.style={font=\sffamily\small\bfseries, text=#1, align=center},
}

\begin{document}
\begin{tikzpicture}

\foreach \name/\title/\note [count=\i] in {
  foto/Foto/{entrada},
  segmentacion/Segmentación/{MediaPipe:\\pelo, piel y fondo},
  marco/Marco de la cabeza/{cráneo superelipse\\y línea del pelo},
  plan/Plan del corte/{se conserva arriba,\\pelo corto a los lados},
  sin_pelo/Quitar el pelo/{relleno pull-push\\solo con fondo seguro},
  corte/Pelo corto/{color propio y\\sombreado de Lambert}} {
  \node[foto] (s\i) at ({2.62 * (\i - 1)}, 0) {\includegraphics[width=2.05cm]{\name.jpg}};
  \node[etapa, above=4pt of s\i] {\title};
  \node[nota, below=3pt of s\i] (n\i) {\note};
}
\foreach \a/\b in {1/2, 2/3, 3/4, 4/5, 5/6} {
  \draw[flecha] (s\a.east) -- (s\b.west);
}

\coordinate (base) at ($(n2.south) + (0, -0.06)$);
\draw[llave=muted] (s2.west |- base) -- (s2.east |- base)
  node[midway, below=9pt, llavetexto=muted] {Modelo};
\draw[llave=blue] (s3.west |- base) -- (s4.east |- base)
  node[midway, below=9pt, llavetexto=blue] {Geometría de la cabeza};
\draw[llave=yellow] (s5.west |- base) -- (s5.east |- base)
  node[midway, below=9pt, llavetexto=yellow] {Reconstrucción};
\draw[llave=red] (s6.west |- base) -- (s6.east |- base)
  node[midway, below=9pt, llavetexto=red] {Render};

\end{tikzpicture}
\end{document}
